% Generated draft: requires mathematical and source review.
\documentclass[11pt]{article}
\usepackage[margin=25mm]{geometry}
\usepackage{fontspec}
\setmainfont{Latin Modern Roman}
\usepackage{amsmath,amssymb,mathtools,unicode-math}
\setmathfont{Latin Modern Math}
\usepackage{enumitem,needspace,xurl,hyperref}
\setlength{\parindent}{0pt}
\setlength{\parskip}{5pt}
\setlength{\emergencystretch}{3em}
\newcommand{\N}{\mathbb{N}}
\newcommand{\Z}{\mathbb{Z}}
\newcommand{\Q}{\mathbb{Q}}
\newcommand{\R}{\mathbb{R}}
\newcommand{\C}{\mathbb{C}}
\newcommand{\F}{\mathbb{F}}
\newcommand{\A}{\mathbb{A}}
\newcommand{\PP}{\mathbb{P}}
\newcommand{\E}{\mathbb{E}}
\newcommand{\eps}{\varepsilon}
\newcommand{\dd}{\,\mathrm{d}}
\begin{document}
\section*{UnsolvedMath ID: 30006222}
\textbf{Source problem code:} catalogue entry
\subsection{Definitions and mathematical statement}
For terminology, in the standard compact-manifold setting, an Anosov flow is a smooth flow \(\varphi_t\) generated by a nowhere-vanishing vector field \(X\) on a compact smooth manifold \(M\) without boundary, with a continuous \(D\varphi_t\)-invariant splitting
\[
TM=\R X\oplus E^s\oplus E^u
\]
and constants \(C,\lambda>0\) such that, for every \(x\in M\) and \(t\geq 0\),
\[
\|D\varphi_t(v)\|\leq Ce^{-\lambda t}\|v\|\quad(v\in E_x^s),
\qquad
\|D\varphi_{-t}(w)\|\leq Ce^{-\lambda t}\|w\|\quad(w\in E_x^u).
\]
This standard definition explains the terminology; it does not resolve the intended geometric hypotheses omitted from the supplied record.

For a finite-dimensional graded cohomology space \(H^\bullet(M;F)\), its determinant line is
\[
\det H^\bullet(M;F)=\bigotimes_q\left(\det H^q(M;F)\right)^{(-1)^q},
\]
where exponent \(-1\) denotes the dual line. In an appropriate analytic-torsion setting, the Ray--Singer metric defines a norm on such a line. The coefficient object \(F\), its required properties, and the geometric data needed for this construction are not specified in the supplied source excerpt.

The source asks three questions concerning analytic torsion and Anosov flows:
\begin{enumerate}
\item Identify a dynamical counterpart of the Ray--Singer metric.
\item Formulate a generalized Fried conjecture relating that dynamical construction to the Ray--Singer metric.
\item Establish the generalized conjecture, described by the supplied registry as a predicted equality between a Ray--Singer norm and a dynamical invariant for general Anosov flows.
\end{enumerate}
The supplied dated assessment attributes a canonical Fried section and a formulation of the generalized conjecture to Jean-Michel Bismut and Shu Shen, and labels the problem partially solved: it reports that the first two questions have been answered, whereas the asserted Ray--Singer norm identity remains conjectural in general despite special cases. This is supplied curation, not an independently verified mathematical conclusion.

The record gives neither the definition of the Fried section nor the formula for the predicted norm identity. It also omits the admissible coefficient systems, geometric assumptions, normalizations, and quantifiers governing that identity. Consequently, the remaining target can only be stated as establishing the cited generalized Fried norm identity under its intended hypotheses; a precise universally quantified equality cannot be recovered from this excerpt without supplying missing mathematical content.
\subsection{Short English statement}
The problem asks for a dynamical analogue of the Ray--Singer metric, a generalized Fried conjecture, and a proof of its predicted norm identity for general Anosov flows. Supplied curation credits Bismut and Shen with the construction and formulation, while reporting the general identity as conjectural. The excerpt omits the identity and its hypotheses, so this statement-only draft requires source consultation and human mathematical review.
\subsection{Sources}
\begin{itemize}
\item Local UnsolvedMath registry, permanent ID \(30006222\), title \emph{Generalized Fried Conjecture for Anosov Flows}, source problem code \texttt{catalogue entry}. The supplied record labels the problem \texttt{partially\_solved} and its statement \texttt{corrected\_verified}; these are registry labels, not independent verification.
\item \emph{Mini-Workshop: Recent Results on Loop Spaces} (2025), Oberwolfach Reports, DOI 10.4171/owr/2025/13, pp.~559--578. Supplied source report and literature reference: \url{https://doi.org/10.4171/owr/2025/13}.
\item Attribution in the supplied original statement: Shu Shen, \emph{Analytic torsion and Anosov flow}, joint work with Jean-Michel Bismut. The registry lists the year 2025 and does not separately identify a proposer.
\item Bismut--Shen, \emph{Anosov vector fields and Fried sections}. Supplied literature reference: \url{https://doi.org/10.1007/s00220-025-05400-8}.
\item Literature assessment checked at \texttt{2026-08-22}: supplied dated curation derived from the Oberwolfach Reports record, not independently verified here. It reports the construction of a canonical Fried section and formulation of a generalized Fried conjecture, with the asserted Ray--Singer norm identity still conjectural in general despite special cases. The record explicitly requests independent verification of time-sensitive status and cited claims.
\end{itemize}
\end{document}
